\section{Conclusion}
I presented MANTA, an end-to-end framework for metadata-only anomaly detection over
encrypted mobile and endpoint traffic. The project combines Android capture, backend-assisted
triage, privacy-aware feature views, privacy-student and federated variants, and a reproducible
public-corpus evaluation pipeline. The results identify a practical boundary: metadata coarsening
can improve release privacy, but it does not provide strong observer privacy against
traffic-analysis attacks.

Answers to the research questions are:
\begin{itemize}[leftmargin=*]
  \item \textbf{RQ1 (metadata-only detection feasibility):} Metadata-only detection is feasible on
  a large public corpus, but the evaluated models remain below deployment-grade quality thresholds.
  The best detector in the March 2026 report matrix is a remote gradient-boosted-tree family variant with F1 \(0.449\),
  PR-AUC \(0.482\), and ROC-AUC \(0.765\).
  \item \textbf{RQ2 (variant comparison):} Remote models outperform the evaluated on-device,
  privacy-student, and federated variants. The privacy-preserving variants are useful for
  measurement, but not yet competitive as detectors.
  \item \textbf{RQ3 (release privacy):} The reduced privacy views lower exact app-ID leakage while
  keeping anomaly utility broadly intact in the March 2026 ablation study. This supports a limited
  release-privacy result.
  \item \textbf{RQ4 (observer privacy):} The encrypted-flow sequence benchmark shows severe
  observer leakage. Release privacy and observer privacy are therefore not interchangeable, and
  feature-level coarsening alone is insufficient for strong traffic-analysis privacy.
\end{itemize}

I claim that MANTA is a reproducible encrypted-traffic IDS framework that separates release privacy
from observer privacy and shows that metadata-only coarsening can reduce exact app
re-identification while remaining insufficient against stronger semantic and passive-observer
inference.

Practically, the project delivers an artifact package: Android app, backend adapter, ML pipeline,
report bundle, manuscript, and reproducibility appendix. Its value is strongest as an empirical
systems study and evaluation framework rather than as a finished privacy-preserving product.

\section{Future Work}
\begin{itemize}[leftmargin=*]
  \item Stronger remote and on-device anomaly models under the same grouped public-corpus protocol.
  \item Transport-level defenses such as padding or timing obfuscation if observer privacy is a hard
  requirement.
  \item Better privacy-student objectives that suppress app-family and behavior leakage more
  strongly while preserving anomaly utility.
  \item More realistic deployment studies on mobile overhead, battery, and long-running user
  acceptance.
  \item Stronger collaborative-learning paths with secure aggregation, client weighting, and
  robustness against poisoning.
\end{itemize}
